\documentclass[12pt,a4paper]{report}
\usepackage[margin=2.54cm]{geometry}
\usepackage{fontspec}
\setmainfont{Times New Roman}
\usepackage{graphicx,xcolor,booktabs,tabularx,longtable,titlesec}
\usepackage[numbers,sort&compress]{natbib}
\usepackage[hidelinks]{hyperref}
\hypersetup{pdftitle={Tìm nạp nền và thông báo đẩy im lặng trên Android}}
\urlstyle{same}
\setlength{\parindent}{0pt}
\setlength{\parskip}{0.4em}
\setlength{\emergencystretch}{2em}
\setlength{\intextsep}{8pt plus 2pt minus 2pt}
\definecolor{diagramink}{HTML}{224563}
\definecolor{diagramfill}{HTML}{EDF4F9}
\newcommand{\diagramnode}[1]{\fcolorbox{diagramink}{diagramfill}{\strut\small #1}}
\renewcommand{\chaptername}{Chương}
\renewcommand{\appendixname}{Phụ lục}
\renewcommand{\contentsname}{Mục lục}
\renewcommand{\listfigurename}{Danh sách hình}
\renewcommand{\listtablename}{Danh sách bảng}
\renewcommand{\bibname}{Tài liệu tham khảo}
\renewcommand{\tablename}{Bảng}
\renewcommand{\figurename}{Hình}
\titleformat{\chapter}[hang]{\normalfont\bfseries\fontsize{16}{20}\selectfont}{\chaptertitlename\ \thechapter.}{0.6em}{}
\titleformat{\section}{\normalfont\bfseries\fontsize{15}{19}\selectfont}{\thesection.}{0.6em}{}
\titleformat{\subsection}{\normalfont\bfseries\fontsize{14}{18}\selectfont}{\thesubsection.}{0.6em}{}
\titlespacing*{\chapter}{0pt}{10pt}{22pt}
\titlespacing*{\section}{0pt}{16pt plus 2pt minus 1pt}{7pt}
\titlespacing*{\subsection}{0pt}{12pt plus 2pt minus 1pt}{5pt}
\setcounter{tocdepth}{1}
\newcommand{\repo}{https://github.com/ThanhTris/Seminar-IE307/blob/db8dc5f}
\newcommand{\source}[2]{\href{\repo/#1}{#2}}

\begin{document}
\hypersetup{pageanchor=false}
\begin{titlepage}
  \centering
  {\bfseries TRƯỜNG ĐẠI HỌC CÔNG NGHỆ THÔNG TIN\par}
  {\bfseries ĐẠI HỌC QUỐC GIA THÀNH PHỐ HỒ CHÍ MINH\par}
  \vspace{1.2cm}
  \includegraphics[height=2.8cm]{assets/uit-logo.png}\par
  \vspace{1.25cm}
  {\fontsize{18}{22}\selectfont\bfseries BÁO CÁO SEMINAR\par}
  \vspace{0.65cm}
  {\fontsize{18}{23}\selectfont\bfseries XỬ LÝ VIỆC TÌM NẠP NỀN VÀ\par
   THÔNG BÁO ĐẨY IM LẶNG TRÊN ANDROID\par}
  \vspace{1.2cm}
  \begin{minipage}{0.9\textwidth}
    \fontsize{13}{18}\selectfont
    \begin{tabularx}{\linewidth}{@{}lX@{}}
      \textbf{Môn học} & IE307.R11 — Công nghệ lập trình đa nền tảng cho ứng dụng di động \\
      \textbf{Giảng viên} & ThS. Phạm Nhật Duy \\
      \textbf{Thành viên} & Trí (trưởng nhóm), Trang, Vinh, Trung, Tâm, Tuấn \\
    \end{tabularx}
  \end{minipage}
  \vfill
  {Thành phố Hồ Chí Minh, tháng 11 năm 2026\par}
\end{titlepage}
\cleardoublepage
\pagenumbering{roman}
\hypersetup{pageanchor=true}

\tableofcontents
\clearpage
\phantomsection
\addcontentsline{toc}{chapter}{\listfigurename}
\listoffigures
\clearpage
\phantomsection
\addcontentsline{toc}{chapter}{\listtablename}
\listoftables
\clearpage
\chapter*{Tóm tắt}
\addcontentsline{toc}{chapter}{Tóm tắt}
\fontsize{13}{18}\selectfont
Báo cáo trình bày ứng dụng bảng tin Android lưu snapshot có phiên bản trong Room và nhận dữ liệu mới khi ở nền. Hệ thống triển khai ba chế độ đồng bộ. PERIODIC dùng WorkManager, PUSH dùng thông báo data-only của Firebase Cloud Messaging và HYBRID kết hợp hai trigger. Cả ba chế độ dùng chung worker lấy snapshot và khâu kiểm phiên bản trước khi ghi cache.

Nhóm đối chiếu kết quả với manifest và log theo từng lượt. Bộ dữ liệu có 23 lượt pilot, gồm 10 lượt hoàn tất hợp lệ theo manifest. OPPO Find X9 ghi nhận năm chu kỳ periodic tự nhiên trong 90 phút và áp dụng snapshot v5. Hai pilot HYBRID ghi nhận ứng dụng lấy được snapshot qua periodic sau khi máy chủ chủ động bỏ gửi push. Các quan sát này mô tả cấu hình pilot đã chạy. Tập chính chưa có lượt đủ điều kiện để so sánh tỷ lệ hoàn thành hoặc độ trễ giữa các cơ chế.

\textbf{Từ khóa.} Android, tải dữ liệu nền, thông báo đẩy im lặng, WorkManager, FCM, Room.
\clearpage

\pagenumbering{arabic}
\fontsize{13}{17}\selectfont

\chapter{Giới thiệu đề tài}
\section{Bài toán và mục tiêu}
Ứng dụng bảng tin cần hiển thị nội dung khi người dùng mở lại, kể cả lúc mạng yếu. Nhóm triển khai bảng tin Android đọc từ Room và kiểm tra khả năng làm mới cache khi người dùng đã rời ứng dụng. Bài toán đặt trọng tâm vào vòng đời ứng dụng và điều kiện chạy nền do Android quản lý.

Demo có máy chủ snapshot, ứng dụng Android và log từ trigger đến ghi cache. Ba chế độ thử là PERIODIC, PUSH và HYBRID. Đồng bộ thủ công dùng để chuẩn bị baseline. Các ca thử kiểm tra target có được apply trong cửa sổ và mạng hoặc Doze ảnh hưởng bước nào. \source{docs/SCOPE.md}{Phạm vi triển khai}.

\section{Cơ sở lý thuyết}
\textbf{Tìm nạp nền.} WorkManager đăng ký lịch với khoảng lặp tối thiểu 15 phút. Android chọn thời điểm thực thi theo lịch và các điều kiện như kết nối mạng \citep{AndroidDefineWork,AndroidDozeStandby}. Lịch đăng ký và worker bắt đầu là hai mốc riêng.

\textbf{Thông báo đẩy im lặng.} FCM data-only mang tín hiệu có thể có bản mới. Sender đặt priority ở mức normal. Callback kiểm tín hiệu rồi tạo công việc cho worker lấy snapshot hiện hành \citep{FirebaseReceiveMessages,FirebaseAndroidPriority}. Firebase chấp nhận lần gửi, Android nhận callback và ứng dụng cập nhật cache là ba mốc riêng.

\textbf{Cache có phiên bản.} Client đọc Room để dựng giao diện. Worker dùng ETag và kiểm version trong transaction trước khi ghi. HYBRID dùng cả lịch periodic và tín hiệu FCM trên cùng đường đồng bộ. \source{docs/MECHANISMS.md}{Mô hình cơ chế}.

\section{Hai nghiên cứu liên quan}
Simon và cộng sự ghi nhận 96.792 trên 100.000 chu kỳ đồng bộ trong mẫu log Mendeley trên iOS không có thay đổi. Nhóm dùng kết quả này để đặt câu hỏi về số GET dẫn đến apply trên bảng tin. Cơ chế giữ lượt lặp của Mendeley gợi ý ca E05-C, khi tín hiệu mới tới lúc worker đang bận \citep{Simon2017DiffsyncEnergy}. \source{research/notes/PAPER_1.md}{Ghi chú SEM-20}.

González-Pérez và cộng sự đánh giá lịch tác vụ trên Android 7--11 bằng AlarmManager, wake lock và thư viện NTD. Cách tác giả đối chiếu lượt dự kiến với lượt quan sát gợi ý ghi cửa sổ, số work start, lượt thiếu và ngoại lệ cho từng cấu hình WorkManager/FCM \citep{GonzalezPerez2022AndroidScheduling}. \source{research/notes/PAPER_2.md}{Ghi chú SEM-21} · \source{research/notes/API_CROSSCHECK.md}{Đối chiếu API SEM-22}.

\chapter{Hiện thực đề tài}
\section{Môi trường thực hiện}
Hồ sơ build Android ghi Windows 11 x64, JBR 25.0.3 (target Java 17), SDK platform 37 và Build Tools 36.0.0. Server yêu cầu Node.js 22+. Dùng Git và Python 3.10+ cho chuỗi xuất và phân tích log. Pilot có emulator Android 17/API 37 và OPPO Find X9 (Android 17, ColorOS V17.0.0). Ba pilot ngày 20/11 dùng app/server \texttt{76fece0}, APK \texttt{814347F\ldots}. Build debug bàn giao có hash \texttt{DBAAD\ldots}. Manifest lưu hash đầy đủ và công cụ build riêng của từng bản. \source{docs/ENVIRONMENT.md}{Hồ sơ môi trường} · \source{release/BUILD_MANIFEST.json}{Provenance build}.

\section{Công cụ và công nghệ sử dụng}
\subsection{Các lớp của hệ thống}
Node.js và TypeScript cung cấp CLI cập nhật, API đọc snapshot và nhật ký sự kiện server. Kotlin và Jetpack Compose hiển thị nội dung lưu trong Room. WorkManager quản lý công việc có ràng buộc mạng, còn FCM chuyển tín hiệu data-only đến callback. Chuỗi xuất log chuẩn hóa sự kiện thành JSONL v2 để analyzer Python đối chiếu với manifest của từng lượt. Hình~\ref{fig:architecture} minh họa cách các thành phần tương tác. \source{docs/SYNC_CORE_API.md}{Hợp đồng sync} · \source{docs/EVENT_SCHEMA.md}{Hợp đồng log}.

\begin{figure}[htbp]
\centering
\renewcommand{\arraystretch}{1.2}
\begin{tabular}{c@{\qquad}c}
\diagramnode{WorkManager đăng ký lịch định kỳ} & \diagramnode{FCM data-only đến callback Android} \\
\small $\downarrow$ & \small $\downarrow$ \\
\multicolumn{2}{c}{\diagramnode{SyncWorker nhận nguồn, mode, run ID và request ID}} \\
\multicolumn{2}{c}{\small $\downarrow$} \\
\multicolumn{2}{c}{\diagramnode{SyncRepository gửi GET tới snapshot server, nhận 200 hoặc 304}} \\
\multicolumn{2}{c}{\small $\downarrow$} \\
\multicolumn{2}{c}{\diagramnode{Room kiểm version và cập nhật cache cùng sự kiện áp dụng}} \\
\multicolumn{2}{c}{\small $\downarrow$} \\
\multicolumn{2}{c}{\diagramnode{Compose đọc bản tin, phiên bản và timeline trong Room}} \\
\end{tabular}
\caption{Kiến trúc đồng bộ bảng tin khi ứng dụng ở nền}
\label{fig:architecture}
\end{figure}

\subsection{Luồng xử lý snapshot}
Server ghi snapshot mới và sự kiện \texttt{server\_updated} cùng một lần commit. Worker gửi \texttt{GET /api/bulletin} với ETag của cache nếu có. Với HTTP 200, repository xác thực phiên bản API và nội dung, rồi so phiên bản nhận được với cache hiện tại trong Room transaction. Nếu bản mới hơn, transaction thay danh sách bản tin, ghi version và ghi \texttt{snapshot\_applied}. Phản hồi cùng phiên bản, phản hồi cũ và HTTP 304 hợp lệ giữ nguyên các hàng bản tin. Version của snapshot được lưu ở metadata để vẫn xác định được bản mới nhất khi danh sách rỗng. \source{docs/API.md}{API snapshot}.

\section{Các bước cài đặt}
\begin{enumerate}
  \item Clone mã nguồn, checkout revision muốn chạy và cài JDK, SDK Android, Node.js và Python theo hồ sơ môi trường.
  \item Trong \texttt{server/}, lần lượt chạy \texttt{npm ci}, \texttt{npm run build} và \texttt{npm test}. Sau khi kiểm thử, chạy server ở một terminal riêng.
  \item Chạy Gradle wrapper trong \texttt{android/}. Dùng \texttt{:app:testDebugUnitTest} để kiểm thử và \texttt{:app:assembleDebug} để build APK. Khi cần APK instrumentation, chạy thêm \texttt{:app:assembleDebugAndroidTest}.
  \item Dùng \texttt{adb devices} để xác nhận thiết bị. Đối chiếu SHA-256 với build manifest trước khi cài bằng \texttt{adb install -r}. Khi thu log bằng instrumentation, dùng cặp app và test APK từ cùng bộ build.
\end{enumerate}
\source{docs/SETUP.md}{Lệnh cài đặt và kiểm tra} · \source{docs/DEMO_RUNBOOK.md}{Runbook kiểm chứng}.

\section{Cấu hình}
Đặt \texttt{HOST=127.0.0.1}, \texttt{PORT=3000} và \texttt{STORE\_DIR} riêng qua biến môi trường của shell. File \texttt{.env.example} ghi các biến cấu hình. Emulator đi qua host gateway \texttt{10.0.2.2:3000}. Thiết bị USB dùng \texttt{127.0.0.1:3000} với \texttt{adb reverse tcp:3000 tcp:3000} trên host server. Kiểm GET từ app bằng đồng bộ thủ công. Đặt run ID, device alias và mode trong app rồi đợi công việc cũ kết thúc trước khi đo. \source{docs/setup/SERVER.md}{Biến môi trường server}.

Đối với PUSH, chuẩn bị cấu hình Firebase đúng package và quyền gửi trên máy thử. Giữ file \texttt{google-services.json}, credential và token ở local, chỉ ghi fingerprint khi kiểm token. Dùng data-only priority normal theo cấu hình protocol. TTL của E04 được giữ cố định để diễn giải cửa sổ reconnect. Periodic dùng khoảng lặp tối thiểu 15 phút và tôn trọng constraint \citep{AndroidDefineWork,FirebaseAndroidPriority}. \source{docs/PUSH.md}{Cấu hình push} · \source{experiments/PROTOCOL.md}{Protocol từng ca}.

\section{Các bước triển khai và thực thi}
\subsection{Chuẩn bị và chạy một lượt}
\begin{enumerate}
  \item Chạy \texttt{npm start}, kiểm \texttt{/health} và tạo snapshot baseline với \texttt{npm run update --} theo hướng dẫn setup.
  \item Cấu hình mode, run ID và alias. Đóng lượt cũ trước khi đổi mode. Đồng bộ thủ công để xác nhận cache baseline, đợi work chuẩn bị kết thúc, xác nhận trạng thái lịch rồi đưa app về Home.
  \item Cập nhật snapshot mục tiêu theo mốc của từng ca ở Bảng~\ref{tab:protocol}. Với PUSH, gửi tín hiệu qua \texttt{npm run send-push} sau khi xác minh cấu hình Firebase. Giữ cùng update ID và version giữa update và lần gửi. \source{docs/setup/FCM_SENDER.md}{Tham số sender}.
  \item Đóng đủ cửa sổ, lưu manifest cùng JSONL rồi hoàn tác cấu hình thử. Chạy \texttt{analyze.py} trên đúng dataset và dùng CSV, summary cùng report để kiểm phase, group, cờ phân loại và checksum.
\end{enumerate}

\subsection{Thiết kế các ca thử}
\begin{table}[htbp]
\centering
\small
\caption{Điều kiện và tiêu chí quan sát của các ca E01–E06}
\label{tab:protocol}
\begin{tabularx}{\textwidth}{@{}l>{\raggedright\arraybackslash}p{4.5cm}>{\raggedright\arraybackslash}X@{}}
\toprule
Ca & Điều kiện thực thi & Dấu hiệu cần đối chiếu \\
\midrule
E01 & PERIODIC tự nhiên, 90 phút & Ít nhất ba chu kỳ là mục tiêu quan sát. Đếm request PERIODIC mới trong cửa sổ và kiểm version cuối. \\
E02 & PUSH, 120 giây từ commit target & Nối update và lần gửi với callback, enqueue, worker và apply. Ghi timeout nếu đủ cửa sổ và target chưa được apply. \\
E03 & 5 phút baseline, 20 phút deep Doze, 10 phút sau thoát & Ghi mốc xác nhận deep IDLE trước khi update. Đối chiếu event theo từng pha. \\
E04 & Offline 2 phút, theo dõi 30 phút PERIODIC hoặc 5 phút PUSH sau nối lại & Update trong lúc offline. Kiểm apply sau reconnect và giữ kết quả theo cửa sổ từng mode. \\
E05 & PUSH, ba tình huống, 15 phút mỗi lượt & Kiểm tín hiệu trùng, hint cũ và tín hiệu mới khi worker bận. Theo dõi version cuối và evidence overlap. \\
E06 & HYBRID, 45 phút từ commit & Server chủ động bỏ gửi. Nối lần apply target với work và fetch có source PERIODIC. \\
\bottomrule
\end{tabularx}
\end{table}

Chuỗi log dùng \texttt{run\_id}, \texttt{device\_alias} và \texttt{signal\_id} để đối chiếu lần gửi và callback. Request ID nối các event của một lượt đồng bộ. Work ID của một lịch periodic có thể lặp qua nhiều chu kỳ. Khi xác nhận thời gian giữa server và Android, cần mẫu clock và sai số phù hợp. \source{experiments/PROTOCOL.md}{Protocol từng ca} · \source{experiments/ANALYSIS.md}{Tái lập phân tích}.

\section{Kết quả đầu ra}
\label{sec:results}
\subsection{Tổng hợp dữ liệu và phạm vi phép đo}
Snapshot SEM-43 tại \texttt{007a1aa} ghi 20 lượt thử SEM-41, gồm tám VALID/SUCCESS sau khi sửa collector và 12 INVALID lịch sử. SEM-45 v0.7.0 tại \texttt{db8dc5f} thêm ba pilot E01 trong PR 35, cho tổng \textbf{23 pilot} thuộc 10 nhóm, với 10 VALID/COMPLETED và 13 INVALID theo manifest. Outcome toàn inventory là 15 SUCCESS, 5 TIMEOUT, 2 FAILURE và 1 UNKNOWN. Hai pilot legacy SEM-29 và các control development được đọc riêng theo SUMMARY.

Kế hoạch main có 33 slot gồm 27 emulator và 6 OPPO. Snapshot bàn giao có 0 run main hợp lệ. Tỷ lệ hoàn thành và thống kê latency main để N/A. Cả ba pilot mới có candidate gate false, cùng build pilot nhưng khác route/thiết bị. \source{experiments/SUMMARY.md}{SUMMARY SEM-43} · \source{experiments/results.csv}{CSV SEM-45} · \source{docs/evidence/SEM-45/CHECKS.md}{CHECKS SEM-45}.

\subsection{Quan sát periodic trên OPPO và emulator}
Bảng~\ref{tab:e01} trình bày ba lượt E01 ngày 20/11. Mục tiêu quan sát của E01 là ít nhất ba chu kỳ tự nhiên trong 90 phút. Điều kiện nguồn điện, tuyến mạng và bản build được giữ cùng từng kết quả.

\begin{table}[htbp]
\centering
\small
\caption{Ba pilot E01 ngày 20/11 theo manifest và Android elapsed}
\label{tab:e01}
\begin{tabularx}{\textwidth}{@{}p{2.7cm}>{\raggedright\arraybackslash}p{3.9cm}>{\raggedright\arraybackslash}X@{}}
\toprule
Lượt & Trạng thái nguồn & Quan sát \\
\midrule
OPPO run-03 & COMPLETED / SUCCESS / VALID & 5 chu kỳ trong 90 phút. Apply v5 tại 930\,021 ms từ đầu cửa sổ Android elapsed. Tuyến ADB reverse, cấp nguồn USB. \\
OPPO run-02 & ABORTED / UNKNOWN / INVALID & App lên tiền cảnh và thiếu mốc mở. Số chu kỳ trong cửa sổ để trống. \\
Emulator 8d5816 & COMPLETED / SUCCESS / VALID & 2 chu kỳ, dưới mục tiêu 3. Tuyến host gateway được phân nhóm riêng với host HTTP của loạt cũ. \\
\bottomrule
\end{tabularx}
\end{table}

Phép đếm mô tả dùng work start PERIODIC, attempt 0, request ID riêng trong khoảng elapsed đã khai báo. Event thiếu boot ID nên cột xác minh theo boot vẫn trống. Mốc 930\,021 ms là thời gian theo Android elapsed tính từ đầu cửa sổ. Clock liên máy chưa kiểm soát nên latency server đến Android để trống. \source{experiments/pilots/run-sem42-oppo-e01-03/manifest.json}{Manifest OPPO run-03} · \source{docs/evidence/SEM-42/CHECKS.md}{CHECKS SEM-42}.

\subsection{Quan sát theo điều kiện thực thi}
Trong E03, hai pilot PERIODIC áp dụng target sau khi thoát force-idle. Trong E04, ba pilot PERIODIC áp dụng target sau reconnect. Cả hai nhóm mô tả hành vi ở các điều kiện đã được tạo trong lượt thử. E06 ghi hai lần server chủ động bỏ gửi và ứng dụng vẫn áp dụng được target từ nguồn PERIODIC. Các lượt E06 dùng APK chưa cấu hình Firebase, nên kết quả mô tả đường dự phòng periodic trong tình huống bỏ gửi chủ động.

E02 FCM đầu cuối và E05 race trên thiết bị ở trạng thái chưa chạy trong bộ dữ liệu. Kiểm thử version guard mô tả hành vi logic của transaction. Bằng chứng cho E05-C còn cần worker đang chạy và tín hiệu mới đã tới trong lúc công việc đó chưa kết thúc. \source{experiments/SUMMARY.md}{Quan sát và giới hạn theo ca} · \source{docs/evidence/SEM-39/CORE_CASES.md}{Kiểm logic cache}.

\section{Minh họa chức năng và thành phần chính}
\label{sec:architecture}
\subsection{Cache và các bảng Room}
Hình~\ref{fig:room-schema} biểu diễn các bảng Room v2 và những cột dùng để theo dõi snapshot. \texttt{bulletin\_items} chứa dữ liệu hiển thị. \texttt{sync\_metadata} giữ khóa cấu hình và trạng thái như \texttt{cachedVersion} và \texttt{lastSyncAtUtc}. \texttt{timeline\_events} giữ mốc thực thi và các trường nối sự kiện theo run, request và worker. Khi apply bản mới, repository thay nội dung cache, cập nhật metadata và ghi event applied trong cùng một transaction. \source{android/app/src/main/java/com/thanhtris/seminarie307/data/local/AppDatabase.kt}{AppDatabase} · \source{android/app/src/main/java/com/thanhtris/seminarie307/sync/SyncRepository.kt}{SyncRepository}.

\begin{figure}[htbp]
\centering
\small
\setlength{\fboxsep}{7pt}
\begin{tabularx}{\textwidth}{@{}X X X@{}}
\fcolorbox{diagramink}{diagramfill}{\begin{minipage}[t][3.5cm][t]{0.86\linewidth}
\textcolor{diagramink}{\texttt{bulletin\_items}}\par\smallskip
\texttt{id} (PK)\\
\texttt{title}, \texttt{body}\\
\texttt{version}\\[5pt]
Nội dung bản tin
\end{minipage}}
&
\fcolorbox{diagramink}{diagramfill}{\begin{minipage}[t][3.5cm][t]{0.86\linewidth}
\textcolor{diagramink}{\texttt{sync\_metadata}}\par\smallskip
\texttt{key} (PK)\\
\texttt{value}\\[5pt]
Phiên bản cache và cấu hình đồng bộ
\end{minipage}}
&
\fcolorbox{diagramink}{diagramfill}{\begin{minipage}[t][3.5cm][t]{0.86\linewidth}
\textcolor{diagramink}{\texttt{timeline\_events}}\par\smallskip
\texttt{eventId} (PK)\\
\texttt{eventType}, \texttt{source}\\
\texttt{runId}, \texttt{requestId}\\
\texttt{workId}, \texttt{version}\\
\texttt{occurredAtUtc}\\[5pt]
Dấu vết xử lý
\end{minipage}}
\end{tabularx}
\caption{Lược đồ ba bảng Room và các cột tiêu biểu}
\label{fig:room-schema}
\end{figure}

\subsection{Đọc giao diện và chuỗi sự kiện}
Màn hình bảng tin hiển thị version cache, mode, run ID và timeline. Trong PERIODIC, theo dõi lần \texttt{work\_started} có source PERIODIC. Với PUSH, đọc chuỗi từ callback đến apply. HYBRID có thể nhận cả hai loại trigger. Khi phân tích E06, kiểm nhật ký máy chủ để xác nhận lượt bỏ gửi và nối lần apply về chu kỳ PERIODIC. Bước mở app sau cửa sổ dùng để đọc cache. \source{android/app/src/main/java/com/thanhtris/seminarie307/ui/BulletinScreen.kt}{Màn hình Compose} · \source{slides/DEMO_SCRIPT.md}{Kịch bản demo}.

\subsection{Timeline của một lượt periodic trên OPPO}
Hình~\ref{fig:oppo-timeline} đặt năm mốc bắt đầu worker của OPPO run-03 trên trục thời gian tính từ đầu cửa sổ. Chu kỳ 1 bắt đầu tại 929\,992 ms. Sau khi GET trả HTTP 200, app ghi \texttt{snapshot\_applied} v5 tại 930\,021 ms với cùng request ID. Bốn chu kỳ sau trả 304 và giữ version 5. Đây là dấu vết của cache đã cập nhật, với năm lần thực thi worker trong cửa sổ. \source{experiments/pilots/run-sem42-oppo-e01-03/events.jsonl}{Timeline canonical run-03}.

\begin{figure}[htbp]
\centering
\setlength{\unitlength}{1mm}
\small
\begin{picture}(153,34)
\put(6,12){\line(1,0){144}}
\multiput(6,11)(24,0){7}{\line(0,1){2}}
\put(6,7){\makebox(0,0){0}}
\put(30,7){\makebox(0,0){15}}
\put(54,7){\makebox(0,0){30}}
\put(78,7){\makebox(0,0){45}}
\put(102,7){\makebox(0,0){60}}
\put(126,7){\makebox(0,0){75}}
\put(150,7){\makebox(0,0){90}}
\put(78,1){\makebox(0,0){Phút tính từ đầu cửa sổ Android elapsed}}
\begingroup\color{diagramink}
\put(30.8,12){\makebox(0,0){\textbullet}} % 15.500 min
\put(56.2,12){\makebox(0,0){\textbullet}} % 31.346 min
\put(80.7,12){\makebox(0,0){\textbullet}} % 46.717 min
\put(108.2,12){\makebox(0,0){\textbullet}} % 63.875 min
\put(132.2,12){\makebox(0,0){\textbullet}} % 78.877 min
\put(30.8,17){\makebox(0,0){\textbf{1}}}
\put(56.2,17){\makebox(0,0){\textbf{2}}}
\put(80.7,17){\makebox(0,0){\textbf{3}}}
\put(108.2,17){\makebox(0,0){\textbf{4}}}
\put(132.2,17){\makebox(0,0){\textbf{5}}}
\put(30.8,24){\makebox(0,0){HTTP 200}}
\put(30.8,30){\makebox(0,0){Apply v5}}
\put(56.2,24){\makebox(0,0){304}}
\put(80.7,24){\makebox(0,0){304}}
\put(108.2,24){\makebox(0,0){304}}
\put(132.2,24){\makebox(0,0){304}}
\endgroup
\end{picture}
\caption{Năm chu kỳ periodic quan sát được trong OPPO run-03}
\label{fig:oppo-timeline}
\end{figure}

\section{Vấn đề gặp phải và cách xử lý}
\begin{itemize}
  \setlength{\itemsep}{2pt}
  \item \textbf{Lịch nền và Doze.} Lịch đăng ký xuất hiện trước work start. Constraint và idle có thể hoãn thực thi \citep{AndroidDefineWork,AndroidDozeStandby}. Audit loạt cũ phát hiện thao tác dump scheduler có khả năng làm thay đổi trạng thái timer. Nhóm thu hồi qualification của các lượt bị ảnh hưởng, giữ outcome cũ và dùng batch chạy lại để mô tả pilot.
  \item \textbf{Token và callback.} Token có thể đổi, normal priority có thể trễ khi Doze, và callback có thời gian xử lý giới hạn \citep{FirebaseReceiveMessages,FirebaseAndroidPriority}. Kiểm cấu hình và token local, chuyển công việc sang worker, đối chiếu update ID, signal ID và kết quả cache.
  \item \textbf{Trùng tín hiệu và response cũ.} Log rejection và correlation giúp truy vết tín hiệu trùng. Version guard trong transaction bảo vệ cache khi response đến khác thứ tự. Khả năng giữ update trong E05-C cần overlap và target cuối từ một lượt riêng.
  \item \textbf{OEM, nguồn điện và clock.} OPPO run-02 bị gián đoạn khi app lên tiền cảnh; run-03 hoàn tất với USB. Các cấu hình khác nhau được giữ riêng. Khi clock hoặc boot thiếu căn cứ, báo count mô tả và giữ duration chưa xác minh ở trạng thái trống.
\end{itemize}
\source{docs/evidence/SEM-41/SCHEDULER_OBSERVER_AUDIT.json}{Audit collector} · \source{docs/evidence/SEM-39/CORE_CASES.md}{Kiểm version guard} · \source{docs/evidence/SEM-42/CHECKS.md}{Evidence thiết bị}.

\section{Video demo}
Bản trình bày demo đi theo ba đoạn. Đoạn PERIODIC đối chiếu lịch đăng ký và chu kỳ nền thật. Đoạn PUSH theo dõi update, callback, worker và snapshot áp dụng. Đoạn HYBRID đối chiếu lượt bỏ gửi phía server và lần worker định kỳ lấy target. Khi dùng playback, hiển thị nhãn và run ID để người xem truy được mốc trong log. \source{docs/DEMO_RUNBOOK.md}{Quy trình demo} · \source{slides/DEMO_SCRIPT.md}{Kịch bản trình chiếu}.

\chapter{Kết luận và khuyến nghị}
\section{Kết luận}
Nhóm đã hiện thực đường đồng bộ chung tới cache và quan sát periodic ở cấu hình pilot emulator và OPPO. E06 có quan sát deliberate skip cùng phục hồi qua periodic. Các lần thử bị hủy, timeout và thiếu evidence được giữ cùng trạng thái nguồn. Những kết quả này mô tả khả năng chạy của cấu hình thử. Phép so latency, năng lượng và khả năng delivery FCM cần mẫu main, clock kiểm soát và điều kiện thiết bị được ghi rõ.

\section{Khuyến nghị}
Periodic phù hợp khi dữ liệu có thể chờ cơ hội chạy nền và cần lượt kiểm tra định kỳ. Push thích hợp cho thay đổi do server phát hiện khi có Firebase và cùng chuỗi evidence nhận đến áp dụng. Hybrid phù hợp khi thiết kế cần cả tín hiệu thay đổi và cơ hội kiểm tra định kỳ, với version guard và quy tắc xử lý trùng rõ ràng. Chọn cơ chế theo yêu cầu độ mới dữ liệu, chi phí GET và điều kiện thiết bị. Sau khi có candidate, ưu tiên chạy E02 đầu cuối và E05-C với overlap quan sát được, đồng thời kiểm cửa sổ và clock trước khi so sánh.

\appendix
\chapter{Bảng phân công công việc}
Phân công theo từng owner từ \source{docs/TEAM_ASSIGNMENTS.md}{TEAM\_ASSIGNMENTS} và \source{tasks/catalog.json}{catalog}. Bảng trình bày đầu việc được giao theo snapshot báo cáo.
\small
\begin{longtable}{@{}p{1.7cm}p{1.55cm}p{11.65cm}@{}}
\toprule
Task & Owner & Đầu việc \\
\midrule
\endhead
\bottomrule
\endlastfoot
SEM-11 & Trí & Chốt phạm vi và môi trường build Android \\
SEM-14 & Trí & Vẽ luồng dữ liệu và hợp đồng sync chung \\
SEM-23 & Trí & Tạo scaffold Android và UI cache/timeline \\
SEM-28 & Trí & Ghép manual sync đầu cuối và chốt hướng dẫn setup \\
SEM-42 & Trí & Chạy đối chiếu OPPO E01/E02 và ghi giới hạn OEM \\
SEM-53 & Trí & Tái lập runbook và quay video dự phòng \\
SEM-15 & Trang & Giải thích cơ chế silent push và giới hạn FCM \\
SEM-22 & Trang & Đối chiếu hai bài nghiên cứu với API push hiện hành \\
SEM-32 & Trang & Triển khai callback FCM và enqueue sync \\
SEM-33 & Trang & Xử lý token refresh và bàn giao PUSH.md \\
SEM-38 & Trang & Xử lý duplicate và version hint của push \\
SEM-47 & Trang & Viết phần silent push và lời dẫn demo 4 phút \\
SEM-12 & Vinh & Lập hồ sơ OPPO và emulator \\
SEM-16 & Vinh & Lập ma trận lifecycle và protocol E01–E06 \\
SEM-24 & Vinh & Chuẩn bị bộ xuất log và run manifest \\
SEM-29 & Vinh & Chạy pilot periodic E01 trên hai thiết bị \\
SEM-34 & Vinh & Chạy pilot push E02 và xác nhận cache nền \\
SEM-41 & Vinh & Chạy tập chính E01–E06 trên emulator \\
SEM-43 & Vinh & Chốt SUMMARY và bộ evidence thực nghiệm \\
SEM-48 & Vinh & Viết nội dung và biểu đồ kết quả cho bài nói \\
SEM-17 & Trung & Giải thích periodic WorkManager và constraints \\
SEM-25 & Trung & Xây Room, SyncRepository và SyncWorker dùng chung \\
SEM-30 & Trung & Triển khai unique periodic work \\
SEM-31 & Trung & Bàn giao chẩn đoán và tài liệu periodic \\
SEM-36 & Trung & Tích hợp ba mode và pilot hybrid \\
SEM-37 & Trung & Chốt bản ứng viên và manifest phục vụ tập chính \\
SEM-39 & Trung & Củng cố version guard, retry và race trong core \\
SEM-52 & Trung & Build APK bàn giao và viết DEMO\_RUNBOOK \\
\pagebreak
SEM-13 & Tâm & Kiểm môi trường Node và quyền Firebase \\
SEM-18 & Tâm & Giải thích server, FCM accepted và correlation \\
SEM-21 & Tâm & Đọc bài nghiên cứu thứ hai và ghi chú nguồn \\
SEM-26 & Tâm & Chốt hợp đồng API và event schema \\
SEM-27 & Tâm & Triển khai server snapshot và CLI cập nhật \\
SEM-35 & Tâm & Làm CLI gửi FCM data-only normal priority \\
SEM-40 & Tâm & Thêm fault injection cố ý bỏ gửi push \\
SEM-44 & Tâm & Kiểm độc lập trigger và phép tính dữ liệu main \\
SEM-49 & Tâm & Viết phần kiến trúc và lập luận nghiên cứu thứ hai \\
SEM-19 & Tuấn & Tổng hợp mô hình cơ chế và câu hỏi phản biện \\
SEM-20 & Tuấn & Đọc bài nghiên cứu thứ nhất và ghi chú nguồn \\
SEM-45 & Tuấn & Viết parser log và sinh results.csv \\
SEM-46 & Tuấn & Viết phần mở bài và bàn giao lời dẫn thao tác demo \\
SEM-50 & Tuấn & Hoàn thiện báo cáo, bibliography và CLAIMS \\
SEM-51 & Tuấn & Ghép deck 15 phút, PDF và speaker notes \\
SEM-54 & Tuấn & Điều phối hai lần diễn tập và ghi thời lượng \\
SEM-55 & Tuấn & Chốt manifest và gói bàn giao seminar \\
\end{longtable}
\normalsize


\clearpage
\addcontentsline{toc}{chapter}{Tài liệu tham khảo}
\bibliographystyle{plainnat}
\bibliography{../../research/references}
\end{document}
